\subsection{Patrón de Reducción Binaria (tipo Josefo)}

\graybold{Preguntas guía:}

\begin{itemize}
\item ¿El problema describe un proceso repetitivo/alternado (cada segundo elemento se "marca") sobre n elementos, con n muy grande (≥10\textasciicircum{}9)?
\item ¿Simular paso a paso sería O(n), demasiado lento, pero el patrón se repite de forma autosimilar al dividir por la mitad?
\item ¿Puedo replantear el problema de tamaño n como un subproblema de tamaño \textasciitilde{}n/2 aplicando la misma regla?
\end{itemize}

\graybold{Utilidad:} Resolver en O(log n) problemas de simulación con patrones alternados sobre secuencias enormes, sin simular cada paso. Es la misma idea del clásico Problema de Josefo: en vez de eliminar uno por uno, se "salta" de una ronda completa a la mitad del problema restante.

\graybold{Complejidad:} O(log n) por consulta (cada ronda aprox. divide n a la mitad).

\graybold{Fundamento teórico:} En cada "ronda" del proceso, la mitad de los elementos son descartados/activados según la paridad de su posición. Reformulando la posición y el tamaño del problema tras cada ronda (n' = n − procesados, i' = nueva posición relativa) se llega a una recurrencia que converge en O(log n) pasos en vez de O(n).

\graybold{Ejercicio aplicado}

Hay N personas en fila y un detector de metales defectuoso que se activa de forma alternada (la 1ª persona no lo activa, la 2ª sí, la 3ª no, ...; quien lo activa va al final de la fila). Dado que el coach está en la posición i, determinar en qué turno (entre quienes no activan el detector) pasará.

\graybold{I/O:} T ≤ 1000 casos, cada uno con solo dos enteros → readline por línea (patrón 1) es suficiente y más simple.

\graybold{Código solución}

\begin{lstlisting}[language=Python]
import sys
input = sys.stdin.buffer.readline
 
def safe_position(n, i):
# turno "seguro" (que no activa el patron) en que pasara la posicion i
    res = 0
    while i % 2 == 0:          # mientras i cae en la mitad "activada"
        res += i // 2         # cuenta los turnos seguros de esta ronda
        n -= i // 2             # reduce el problema a la mitad restante
        i = n                   # la posicion i se convierte en el nuevo n
    res += (i + 1) // 2        # suma los turnos seguros de la ronda final
    return res
 
def solve():
    t = int(input())
    out = []
    for _ in range(t):
        n, i = map(int, input().split())
        out.append(str(safe_position(n, i)))
    sys.stdout.write("\n".join(out))
 
solve()
\end{lstlisting}